\begin{statementbox}
\textbf{条件与记号。} $\triangle ABC$ 为非退化三角形，记 $AB=p>0,AC=q>0,BC=a>0$。
点 $D$ 在线段 $BC$ 的\textbf{内部}，即顺序为 $B,D,C$，所有线段符号均表示普通正长度。

\textbf{正向定理与逆定理：}
\[
 \boxed{\angle BAD=\angle DAC
 \quad\Longleftrightarrow\quad
 \frac{BD}{DC}=\frac{AB}{AC}=\frac pq}.
\]
其中由左到右为内角平分线定理，由右到左为其逆定理；两者都保留 $D$ 在 $BC$ 内部的前提。
也可写成 $BD\cdot AC=DC\cdot AB$。

\textbf{按邻边权重分配整段：} 因 $BD+DC=a$，有
\[
 \boxed{BD=\frac{ap}{p+q},\qquad DC=\frac{aq}{p+q}}.
\]
若取 $B=0,C=a$，分点坐标就是 $x_D=ap/(p+q)$，严格落在 $(0,a)$ 内。
较长的邻边对应较长的分段：$p>q$ 时 $BD>DC$，点 $D$ 更靠近 $C$；
$p<q$ 时方向相反。

\textbf{何时也是中线或高？} 在 $AD$ 已是内角平分线的条件下，
\[
 AB=AC\quad\Longleftrightarrow\quad BD=DC
 \quad\Longleftrightarrow\quad AD\perp BC.
\]
一般三角形中，角平分线、中线和高不能混用。
只有两邻边相等时，三条线才重合；本条不需要排除等腰三角形。
\end{statementbox}
